
\chapter{Комплаенс и AML}\label{ch:compliance-aml}
Комплаенс в платёжном бизнесе начинается с обычного вопроса: кого
именно подключаем, что уже известно о клиенте
или ТСП и в какой момент нужно остановить операцию. Для
платёжного продукта это рабочий механизм, который срабатывает внутри
реального потока операций.

\section{115-ФЗ и платёжный бизнес}\label{sec:compliance-115fz}

115-ФЗ, Федеральный закон \enquote{О противодействии
  легализации (отмыванию) доходов, полученных преступным путём, и
финансированию терроризма}, задаёт базовую рамку AML-контроля в
России~\cite{fz-115}. Для платёжного оператора это значит:
пока клиент ещё не прошёл проверку, с его операцией
нельзя обращаться как с полностью доверенной. Закон раскладывает
эту задачу на три обязанности: знать клиента, наблюдать операции и
при формальных основаниях блокировать средства. AML встроен
в сам продукт, а не вынесен во внешний формальный контур.

Первая обязанность касается идентификации. Закон говорит об
идентификации и об упрощённой идентификации, а подзаконные акты Банка
России задают состав сведений и документов, которые надо собрать по
клиенту, представителю, выгодоприобретателю и бенефициарному
владельцу~\cite{cbr-499p}. Для оператора это значит, что на
одних регистрационных данных далеко не уедешь: где-то достаточно
минимального набора сведений, а где-то нужны документы, внешние
реестры и дополнительные вопросы. Конкретный объём проверки зависит
от типа клиента, операции и уровня риска, поэтому в печатной версии
лучше не привязываться к быстро меняющимся порогам и частным
исключениям.

Вторая обязанность связана с мониторингом. Оператор должен замечать
операции, которые подпадают под критерии обязательного контроля,
и отправлять сведения в Росфинмониторинг. Но этим задача не
ограничивается: надо ловить и те операции, которые формально ещё не
описаны в законе, но выбиваются из нормального профиля клиента по
поведению или по контексту. Система требует регламента, потока данных, маршрута проверки и
понятного решения по результату. Если контроль настроен формально,
он либо ничего не увидит, либо захлебнётся в ложных срабатываниях.

Третья обязанность связана с замораживанием (блокированием) денежных
средств и иного имущества. Если клиент или бенефициар попадает в
перечни лиц, связанных с терроризмом или экстремизмом, организация
обязана применить меры по замораживанию (блокированию) и исполнить
предусмотренный законом порядок информирования~\cite{fz-115}. Скорость здесь не отменяет осторожности: для AML промедление иногда опаснее жёсткости.

Связь с лицензированием прямая. Банк России при надзоре за
оператором платёжной инфраструктуры проверяет формальные
правила и то, есть ли в компании рабочая AML-система,
обученные сотрудники и понятный маршрут обработки кейсов. 115-ФЗ
задаёт конкретную операционную обязанность. AML-контур начинается
раньше первой спорной операции. Он определяет, кого продукт вообще
готов впустить и на каких условиях готов продолжать работу.

\section{KYC при подключении физлиц}\label{sec:compliance-kyc}

Для физлица KYC (Know Your Customer) начинается с вопроса о степени
доверия: можно ли считать этого человека
тем, за кого он себя выдаёт, и какой риск берём на себя, открывая
ему следующий уровень сервиса. В платёжном приложении после регистрации оператор решает, какие лимиты
оставить сразу, а какие открывать только после дополнительной
проверки. KYC работает как лестница уровней доверия.

На практике эта лестница идёт от минимального набора сведений к
документальной сверке и дополнительным проверкам там, где этого
требует закон или риск-профиль~\cite{cbr-499p}. Чем меньше
сведений собрано на входе, тем уже набор доступных операций и
лимитов; Банк России прямо разъясняет это на примере переводов и
электронных кошельков по упрощённой идентификации~\cite{cbr-banking-faq-identification}. Автоматизация делает вход
быстрее, однако не отменяет базовой обязанности проверять качество
данных: ошибка на входе редко остаётся локальной. Поэтому KYC почти
всегда балансирует между удобством и точностью.

Набор проверок: сверка документов
по государственным источникам, проверка по перечням
Росфинмониторинга и иным санкционным спискам, а также выявление
политически значимых лиц (ПЭП)~\cite{cbr-499p,fatf-pep-guidance}.

\definitionnote{ПЭП}{Человек, наделённый или ранее наделённый
  значимой публичной функцией; в AML-контуре отдельно учитывают также
  членов семьи и близких соратников. Для такой категории применяют
усиленную проверку из-за повышенного коррупционного риска.}

Для неидентифицированного пользователя электронных денежных
средств 161-ФЗ (статья~10) ограничивает остаток ЭДС и месячный
оборот: остаток ЭДС не более 15\,000~руб. в любой момент
и оборот не более 40\,000~руб. в месяц без прохождения
упрощённой идентификации. После упрощённой идентификации
лимит остатка ЭДС поднимается до 60\,000~руб., а месячный
оборот -- до 200\,000~руб.~\cite{fz-161}. Для полной
(персонифицированной) идентификации физлица-резидента РФ
лимит остатка установлен законом на уровне 600\,000~руб.;
дополнительные параметры могут определяться договором
с оператором. Эти пороги задают архитектуру уровней: продукт
должен хранить текущий статус идентификации клиента и применять
соответствующий лимитный контур до каждой операции.

KYC для физлица устроен как постепенное наращивание доверия. Чем
чувствительнее сценарий
и выше риск, тем меньше система может полагаться на самодекларацию
и тем больше ей нужны внешние подтверждения. У ТСП задача
сложнее ещё на один слой: там приходится проверять человека и сам
бизнес как действующую конструкцию.

\section{KYB при подключении ТСП}\label{sec:compliance-kyb}

Проверка компании отличается от проверки физлица прежде всего
тем, что у юридического лица нужно оценить личность и деловую
конструкцию: структуру собственности, реальный вид
деятельности, возвраты, спорные операции и риск использования
бизнеса как транзитной оболочки. Платформа, которая подключает
нового продавца, должна понять, кто за ней стоит и как она
действительно зарабатывает. Поэтому KYB
(Know Your Business) собирает регистрационные данные, владельцев
и платёжное поведение в единый профиль риска.

Базовый пакет документов включает выписку из реестра, данные
директора и учредителей, сведения о бенефициарном владельце
(ultimate beneficial owner, UBO), банковские реквизиты и описание
бизнеса. Если структура владения не сходится, а сайт не похож на
заявленный вид деятельности, вопросы должны возникнуть сразу.
Проверка строится на сопоставлении официальных данных с тем, как
бизнес выглядит в реальности. Но формально чистая компания вполне
может оказаться рискованной на практике. Именно поэтому KYB нельзя
заканчивать папкой с документами: после подключения его приходится
постоянно сверять с живым поведением бизнеса. В логике
идентификации по 115-ФЗ и \mbox{499-П} это прямо предусмотрено:
кредитная организация должна установить данные самого
клиента и сведения о бенефициарном владельце~\cite{cbr-499p}.

MCC сам по себе не \enquote{плохой}; схемные правила требуют
корректно классифицировать ТСП и передавать правильный код
категории; для отдельных сегментов и для категорий, которые Visa
относит к high-integrity risk merchants, действуют отдельные
ограничения и требования мониторинга~\cite{visa-core-rules}.
Подмена MCC означает попытку спрятать реальный профиль бизнеса под
более мягкой категорией. Если заявленная категория не сходится с
фактической деятельностью, платёжный поток лучше останавливать до
выяснения.

Модель платёжного фасилитатора (payment facilitator, PayFac)
меняет место ответственности, но не смысл проверки. Visa прямо
возлагает на эквайера ответственность за Payment Facilitator и его
Sponsored Merchants: до начала обработки нужно подтвердить
регистрацию PayFac у Visa, присвоить уникальные идентификаторы
самому PayFac и его Sponsored Merchants, получать отчётность по их
активности и следить за корректным MCC~\cite{visa-core-rules}. Поэтому платформа может подключать
субмерчантов быстрее, но расплачивается за это собственными
процедурами KYB и более плотным мониторингом. Иначе скорость
подключения превращается в скорость накопления риска.

KYB проверяет, совпадает ли юридический фасад бизнеса с его реальным
платёжным поведением.

\practicenote{%
  Если вы проверяете ТСП, смотрите на документы и на то, как
  бизнес выглядит на сайте, какие у него товары, как устроены возвраты
  и нет ли противоречия между заявленным MCC и фактической
  деятельностью. В KYB ошибка в одну сторону стоит
  регуляторного риска, а ошибка в другую оборачивается лишним отказом
  нормальному бизнесу.

}
\begin{table}[tbp]
  \begingroup
  \scriptsize
  \setlength{\tabcolsep}{4pt}
  \renewcommand{\arraystretch}{1.2}
  \makebox[\textwidth][l]{%
    \begin{minipage}{\fullwidth}
      \begin{tabularx}{\fullwidth}{
          @{}>{\RaggedRight\arraybackslash\bfseries}p{0.18\fullwidth}
          Y
          Y
          Y@{}
        }
        \toprule
        \rowcolor{Panel}
        Критерий &
        \cellcolor{Good!16}\textbf{Низкий риск} &
        \cellcolor{Warn!18}\textbf{Средний риск} &
        \cellcolor{Bad!16}\textbf{Высокий риск} \\
        \midrule
        \rowcolor{Soft}
        Примеры MCC &
        retail, SaaS &
        travel &
        crypto exchange,\newline gambling, adult \\
        Документы &
        стандартный пакет &
        расширенный пакет\newline + лицензии &
        расширенный пакет\newline + финансовая отчётность \\
        \rowcolor{Soft}
        Мониторинг после подключения &
        квартальный &
        ежемесячный &
        еженедельный \\
        Rolling reserve &
        нет &
        5--10\% &
        10--20\% \\
        \bottomrule
      \end{tabularx}
      \par\smallskip
      \fcolorbox{Ink!18}{Soft}{%
        \parbox{\dimexpr\fullwidth-2\fboxsep-2\fboxrule\relax}{%
          \textcolor{Muted}{Цвет кодирует не \enquote{хорошо/плохо}, а глубину due
            diligence: чем выше риск-профиль
            ТСП, тем шире пакет документов, плотнее мониторинг
          после подключения и вероятнее rolling reserve.}%
        }%
      }
    \end{minipage}%
  }
  \endgroup
  \caption{Чек-лист KYB по категориям риска.}\label{fig:compliance-kyb-risk}
\end{table}

\section{Санкционный скрининг}\label{sec:compliance-sanctions}

Санкционный скрининг обязателен в тех точках, где
бизнес может вступить в отношения с запрещённым лицом или
юрисдикцией. Если клиент, бенефициар или контрагент попадает в
один из применимых перечней, операцию нужно остановить до исполнения.
Проверка строится на сопоставлении имён, документов и реквизитов с
актуальными перечнями. Трудность в том, что списки длинные, имена
повторяются, а совпадение ещё не означает реальный риск.

Для российского оператора первичны три национальных
перечня. Во-первых, перечень организаций и физических лиц,
в отношении которых имеются сведения об их причастности
к экстремистской деятельности или терроризму, который ведёт
Росфинмониторинг; меры по замораживанию средств таких лиц
применяются по подпункту~6 пункта~1 статьи~7 115-ФЗ
(операции с лицами из перечня дополнительно подлежат
обязательному контролю по статье~6). Во-вторых, перечень
организаций и физических лиц, в отношении которых имеются
сведения об их причастности к распространению оружия
массового уничтожения. В-третьих, решения Межведомственной
комиссии (МВК) о замораживании (блокировании) денежных
средств или иного имущества~\cite{rosfinmonitoring-lists,fz-115,cbr-mvk}.\footnote{Точные
наименования перечней, формат публикации и периодичность
обновления сверяются по сайту Росфинмониторинга и по
действующей редакции 115-ФЗ.} Перечни ООН, ЕС и OFAC SDN
для российского оператора имеют другой статус: их использование
определяется собственной риск-политикой, требованиями
банков-корреспондентов и санкционным режимом конкретной
транзакции, но они не подменяют национальные
перечни~\cite{ofac-sanctions-list-service,eu-sanctions-info,unsc-consolidated-list}.

Помимо санкционного скрининга, 115-ФЗ разводит два режима
контроля операций. Обязательный контроль (статья~6) применяется
к закрытому перечню операций -- по сумме (1\,000\,000~руб.\
для денежных операций с наличными и ряда других) и по виду
(операции с лицами из перечня экстремистов, операции
с недвижимостью свыше порога и т.\,д.). Факультативный контроль
(статья~7) обязывает оператора выявлять необычные
и подозрительные операции на основании собственных правил
внутреннего контроля. Сведения об операциях, подлежащих
обязательному контролю, и о подозрительных операциях
направляются в Росфинмониторинг через личный кабинет в формате
формализованного электронного сообщения (ОЭС) в установленные
115-ФЗ и нормативными актами Банка России и
Росфинмониторинга сроки~\cite{fz-115,cbr-860p,rosfinmonitoring-oes}.\footnote{Состав
показателей и форматы ОЭС регулируются приказом Росфинмониторинга
№ 18 от 08.02.2022 и связанной с ним нормативной базой;
требования к правилам внутреннего контроля кредитных организаций
с 13.08.2025 устанавливает Положение Банка России 860-П
(заменило 375-П). Точные реквизиты сверяются перед публикацией.}

На практике оператор ставит проверку по перечням как минимум
на подключении, при операциях, попадающих в риск-профиль,
при обновлении данных клиента и при обновлении самих
перечней. Перечни Росфинмониторинга обновляются регулярно;
после публикации обновлённого перечня оператор обязан в
установленный срок не позднее одного рабочего дня
прогнать обновление и применить меры к клиентам, попавшим
в перечень~\cite{rosfinmonitoring-lists,fz-115}.\footnote{Срок
применения мер к клиентам после обновления перечня -- предмет
последовательных правок 115-ФЗ; перед публикацией сверить
по действующей редакции.}

Скрининг одной транзакции проходит так.
Авторизационный запрос поступает на хост эмитента. Ещё до
проверки лимитов модуль скрининга извлекает имя держателя
из CMS по PAN, нормализует его (транслитерация, удаление
диакритики, разбиение на n-граммы) и сравнивает с индексом
санкционных перечней в оперативной памяти. Если нечёткое
совпадение превышает порог, транзакция уходит в очередь
ручной проверки; оператор видит контекст (страна, сумма,
история клиента) и принимает решение: снять блок или
эскалировать. Если совпадение ниже порога, транзакция
продолжает обычный путь авторизации; добавленная латентность
ограничена временем lookup в индексе.\footnote{Конкретные
числовые пороги (доля совпадения по Jaro--Winkler и пр.)
и время разбора кейса в книге не приводятся: они зависят
от настройки конкретного оператора и не зафиксированы в
нормативных актах.}

Скрининг в онлайне тормозит авторизацию, и бизнес неизбежно ищет
компромисс между скоростью и полнотой. Поэтому рабочая система
обычно строится на индексах в оперативной памяти и на чёткой
очереди эскалации: ручная проверка каждого совпадения здесь просто
не масштабируется. Распространённые имена и
совпадения между кириллицей и латиницей порождают массу шума;
задача -- находить подозрительные совпадения и быстро разбирать ложные.

Ложноположительные совпадения требуют отдельного
маршрута разбора. Типовой сценарий: система помечает
операцию как требующую проверки, сотрудник смотрит контекст, после
чего либо снимает блок, либо отправляет случай на эскалацию. Такой
маршрут должен быть формализован, иначе итог начнёт зависеть от
того, кто именно дежурит на смене. Если порог совпадения слишком
мягкий, команда утонет в шуме; если слишком жёсткий, пропустит
риск. Логика сопоставления и пороги срабатывания должны
настраиваться исходя из собственной оценки риска оператора,
зафиксированной в правилах внутреннего контроля по 115-ФЗ.

\section{Мониторинг ТСП после подключения}\label{sec:compliance-merchant-monitoring}

Риск меняется после подключения, поэтому мониторинг не
заканчивается на подключении. Схемы типа bust-out и попытки
использовать ТСП для отмывания проявляются именно здесь:
пока профиль выглядит спокойно, система расслабляется, а потом
видит проблему слишком поздно. Поэтому защита должна отслеживать
изменения объёма, географии и доли возвратов.

Поведенческие сигналы для расследования:
резкий рост оборота, нетипичная география клиентов, высокий процент
возвратов и несоответствие MCC реальной деятельности. Без контекста даже чистая статистика способна
ввести в заблуждение.

В промышленных системах это обычно выглядит как слой поверх
авторизационной базы: агрегаты по ТСП, скользящие окна,
поиск выбросов и очередь кейсов для аналитиков. Если у ТСП
внезапно растёт объём операций, команда должна увидеть это почти
сразу, а не ждать конца месяца. В европейском регионе Visa прямо
требует от эквайера сохранять ежедневные данные по обороту, среднему
чеку, числу операций, времени до расчёта и числу диспутов, а затем
сравнивать текущие показатели с нормальным профилем активности как
минимум ежедневно~\cite{visa-core-rules}. Такая система переводит наблюдение в действие; чем больше
сигналов, тем важнее уметь отсекать шум.
Когда мониторинг после подключения не работает, частный сбой быстро
превращается в регуляторный кризис.

\section{Кейс QIWI: уроки для операционного риска}\label{sec:compliance-qiwi}

21 февраля 2024 года Банк России отозвал лицензию у КИВИ Банка~\cite{cbr-qiwi-license-2024}. В официальном пресс-релизе
регулятор указал на систематические нарушения требований
законодательства по ПОД/ФТ, вовлечённость банка в высокорисковые
операции, включая переводы в пользу криптообменников и нелегальных
онлайн-казино, а также на случаи открытия QIWI-кошельков с
использованием персональных данных граждан без их ведома.

Инженерный разрез нарушений выглядит так. 115-ФЗ требует от банков
идентифицировать клиентов и проводить контроль операций; открытие
кошельков без ведома граждан означает, что процедура идентификации
либо не выполнялась, либо не проверялась системой контроля.
Переводы в пользу криптообменников и казино становятся проблемой
тогда, когда transaction monitoring не маркирует их как высокорисковые
или маркирует, но сигнал не попадает на ручной разбор. Оба сбоя --
технические: отсутствие автоматического триггера для enhanced due
diligence и неполный audit trail по спорным транзакциям.

Этот эпизод показывает, что регуляторный риск копится тихо. Если
компания полностью зависит от одного провайдера, она неизбежно
наследует и его слабые места. В момент отзыва лицензии бизнес теряет
рабочий платёжный канал. Поэтому диверсификация провайдеров остаётся
базовой операционной гигиеной.

Из подобных кейсов следуют два вывода. Регуляторные предписания надо
читать как предупреждение, а не как формальность: рост объёма без
роста контроля создаёт слепые пятна. Если один платёжный провайдер
становится единственной точкой отказа, бизнесу нужен заранее
подготовленный запасной маршрут.

\importantnote{%
  Для бизнеса с зависимостью от одного платёжного
  провайдера нужен заранее проверенный план непрерывности:
  альтернативный провайдер, резервные счета и маршрут переключения.
  Отзыв лицензии нельзя отлаживать постфактум.

}
\section{FATF и ЕАГ в СНГ-контексте}\label{sec:compliance-fatf}

FATF -- международный язык риска для платёжного бизнеса:
рекомендации, взаимные оценки и публичные списки юрисдикций
под усиленным наблюдением~\cite{fatf-recs,fatf-increased-monitoring}.
Для российского оператора важны две оговорки. Первая:
24 февраля 2023~года членство Российской Федерации в FATF
было приостановлено~\cite{fatf-russia-suspension-2023};
формально это означает, что Россия не участвует
в выработке решений FATF, не направляет своих представителей
в её рабочие группы и не проходит очередные взаимные оценки
в стандартном порядке. Вторая: на пространстве СНГ функцию
региональной FATF-style body исполняет Евразийская группа
по противодействию легализации преступных доходов и финансированию
терроризма (ЕАГ); именно ЕАГ проводит взаимные оценки государств
региона и публикует методологические документы, на которые
опираются национальные регуляторы стран
СНГ~\cite{eag-overview,eag-mutual-evaluations}.\footnote{Точная
дата приостановки членства РФ в FATF и текущий статус
взаимодействия (приостановлено vs.\ исключено) сверяются
по сайту FATF и официальным сообщениям Росфинмониторинга
перед публикацией.}

Даже после приостановки членства FATF-рамка продолжает
влиять на платёжный бизнес в России косвенно. Банки-корреспонденты
и платёжные посредники в третьих юрисдикциях по-прежнему
смотрят на риск через рекомендации FATF и публичные списки.
Крупные банки практикуют политику снижения контрагентского
риска (de-risking) и сокращают или полностью прекращают
отношения с целыми классами клиентов или
юрисдикциями~\cite{fatf-correspondent-banking}; для
российской стороны это выражается прежде всего в трудностях
открытия и поддержания корреспондентских счетов в валютах
недружественных стран. На практике
это означает дополнительные проверки и ограничение
доступности канала.

\practicenote{%
  Перед выходом в новую юрисдикцию проверяются три рамки:
  оценка ЕАГ (для стран СНГ) или FATF (для остальных),
  последняя публичная взаимная оценка страны и местная
  практика AML. Если один из трёх пунктов вызывает сомнение,
  заранее закладывается запас по времени и каналам расчёта.

}
